\documentclass[11pt,a4paper,reqno]{amsart}

\usepackage[T1]{fontenc}
\usepackage[utf8]{inputenc}
\usepackage[ngerman]{babel}
\usepackage{amsmath,amssymb,amsthm,amsfonts,mathtools}
\usepackage{geometry}
\usepackage{microtype}
\usepackage{xcolor}
\usepackage{tikz-cd}
\usetikzlibrary{babel}
\usepackage{enumitem}
\usepackage{booktabs}
\usepackage{array}
\usepackage{hyperref}

\geometry{tmargin=2.5cm,bmargin=2.5cm,lmargin=2.5cm,rmargin=2.5cm}
\hypersetup{
  colorlinks=true,
  linkcolor=blue!50!black,
  citecolor=blue!50!black,
  urlcolor=blue!50!black,
  pdftitle={Diskrete EABC-Sektorlabels, V4-Fourier-Zerlegung und Hardy-Kompatibilitaet},
  pdfauthor={Thomas Hoffbauer}
}
\setlength{\emergencystretch}{2em}
\raggedbottom

% Theorem-Umgebungen
\newtheorem{theorem}{Theorem}[section]
\newtheorem{lemma}[theorem]{Lemma}
\newtheorem{proposition}[theorem]{Proposition}
\newtheorem{corollary}[theorem]{Korollar}
\theoremstyle{definition}
\newtheorem{definition}[theorem]{Definition}
\newtheorem{hypothesis}[theorem]{Hypothese}
\newtheorem{example}[theorem]{Beispiel}
\theoremstyle{remark}
\newtheorem{remark}[theorem]{Bemerkung}

% Abkürzungen
\newcommand{\M}{\mathcal{M}}
\newcommand{\N}{\mathbb{N}}
\newcommand{\R}{\mathbb{R}}
\newcommand{\ZZ}{\mathbb{Z}}
\newcommand{\CC}{\mathbb{C}}
\newcommand{\FF}{\mathbb{F}}
\newcommand{\Pur}{\operatorname{Pur}}
\newcommand{\Supp}{\operatorname{supp}}
\newcommand{\im}{\operatorname{im}}
\newcommand{\Aut}{\operatorname{Aut}}
\newcommand{\id}{\operatorname{id}}
\newcommand{\End}{\operatorname{End}}
\newcommand{\SEABC}{\mathcal{S}_{EABC}}
\newcommand{\Oh}{\mathcal{O}}
\newcommand{\StatusH}{\textnormal{\textbf{[H]}}}
\newcommand{\StatusA}{\textnormal{\textbf{[A]}}}
\newcommand{\StatusC}{\textnormal{\textbf{[C]}}}
\newcommand{\StatusORQ}{\textnormal{\textbf{[ORQ]}}}

\title[EABC-Sektoren, Fourier-Zerlegung und Hardy]{Diskrete EABC-Sektorlabels, $V_4$-Fourier-Zerlegung und Hardy-Kompatibilität}
\author{Thomas Hoffbauer}
\address{Computer Institut Bamberg, Deutschland}
\email{research@cib.de}
\date{5. August 2026}

\begin{document}

\begin{abstract}
We analyze the precise interface between the discrete EABC prime-classification framework and Hardy's operational reconstruction of quantum theory. Within Hardy's axiomatic framework, the subspace and composite-system axioms imply $K=N^r$ for a positive integer $r$; the continuity axiom excludes $r=1$, and the simplicity axiom then selects $r=2$. None of these conclusions is derived from EABC arithmetic itself. We prove a no-go result: a nontrivial discrete EABC signature space cannot satisfy Hardy's continuity axiom because every continuous path into a discrete space is constant. EABC signatures can instead serve as fiber indices or component labels for a continuous state model. We formulate a conditional operator lift with quotient $V_4$ and prove the finite Fourier decomposition associated with any complex linear representation of $V_4$. If the action is by algebra automorphisms, the Fourier sectors inherit the EABC multiplication law. Formally established statements, source-dependent Hardy results, conditional model assumptions, and open questions are separated by the tags \textbf{[A]}, \textbf{[H]}, \textbf{[C]}, and \textbf{[ORQ]}.
\end{abstract}

\maketitle

\section{Einleitung und Geltungsbereich}

Die EABC-Arithmetik ordnet natürlichen Zahlen drei diskrete Daten zu: einen Primfamilienträger, einen Paritätsvektor und eine multiplikative Restklassensignatur. Parallel dazu erlaubt die Fasersprache, mehrere Modellrepräsentanten über demselben skalaren Wert zusammenzufassen. Die vorliegende Arbeit untersucht, ob und in welchem präzisen Sinn diese diskreten Strukturen mit Hardys Kontinuitätsaxiom für reine Zustände vereinbar sind.

Die Zielsetzung ist bewusst begrenzt. Es wird weder eine Herleitung der Quantenmechanik aus der EABC-Arithmetik noch eine physikalische Interpretation der Primfamilien behauptet. Bewiesen werden ein topologisches No-Go-Resultat, eine algebraische Existenzkonstruktion für diskret-kontinuierliche Operatorerweiterungen und eine bedingte $V_4$-Fourier-Zerlegung. Die Identifikation dieser Konstruktionen mit einem operationalen Quantensystem bleibt eine zusätzliche Modellhypothese.

Wir verwenden folgende epistemische Kennzeichnung:
\begin{itemize}[leftmargin=2.2em]
  \item \StatusH{} bezeichnet Aussagen, die innerhalb von Hardys vollständigem Axiomensystem gelten und aus seiner Arbeit übernommen werden;
  \item \StatusA{} bezeichnet in dieser Note formal bewiesene mathematische Aussagen;
  \item \StatusC{} bezeichnet zusätzliche, mathematisch konsistente Modellannahmen;
  \item \StatusORQ{} bezeichnet offene Forschungsfragen.
\end{itemize}

\section{Hardys Parameter \texorpdfstring{$N$ und $K$}{N und K}}

Hardy charakterisiert einen Systemtyp durch zwei positive ganze Zahlen \cite{Hardy2001}:
\begin{itemize}[leftmargin=2.2em]
  \item $N$ ist die maximale Anzahl von Zuständen, die in einer Einzelmessung zuverlässig voneinander unterschieden werden können;
  \item $K$ ist die minimale Anzahl von Wahrscheinlichkeitsmessungen beziehungsweise reellen Parametern, die den Zustand vollständig bestimmen.
\end{itemize}

Hardys Axiom 4 fordert für ein zusammengesetztes System $AB$
\begin{equation}
N_{AB}=N_A N_B,
\qquad
K_{AB}=K_A K_B.
\end{equation}
Axiom 5 verlangt eine kontinuierliche reversible Transformation zwischen je zwei reinen Zuständen desselben Systems \cite{Hardy2001}.

\begin{proposition}[Potenzform in Hardys Axiomensystem \StatusH]
\label{prop:hardy-power}
Innerhalb von Hardys Axiomensystem folgt
\begin{equation}
K(N)=N^r
\qquad\text{für ein festes }r\in\N_{>0}.
\end{equation}
Dabei liefert das Subraumaxiom die strenge Monotonie von $K$, während das Axiom für zusammengesetzte Systeme die vollständige Multiplikativität
\begin{equation}
K(mn)=K(m)K(n)
\end{equation}
erzwingt. Hardy verwendet anschließend das zahlentheoretische Resultat, dass eine streng wachsende vollständig multiplikative Funktion dieser Form eine Potenzfunktion ist; die Ganzzahligkeit von $K$ führt in seiner Argumentation zu einem positiven ganzzahligen Exponenten.
\end{proposition}

\begin{remark}
Proposition~\ref{prop:hardy-power} ist keine Folgerung aus der EABC-Signatur. Sie gilt erst unter Hardys operationalen Axiomen, insbesondere unter dem Subraum- und Kompositionsaxiom. Die EABC-Multiplikation liefert für sich allein weder eine Zustandsdimension $N$ noch eine Parameterzahl $K$.
\end{remark}

\begin{proposition}[Ausschluss- und Selektionsschritt nach Hardy \StatusH]
\label{prop:hardy-selection}
Hardys Axiom 5 schließt den klassischen Zweig $r=1$ aus. Das Einfachheitsaxiom wählt anschließend den kleinsten noch zulässigen Exponenten, also $r=2$.
\end{proposition}

\begin{remark}
Die logische Reihenfolge lautet
\begin{equation}
\boxed{
\text{Kontinuität verwirft }r=1;
\qquad
\text{Einfachheit wählt danach }r=2.
}
\end{equation}
Axiom 5 allein beweist daher nicht $K=N^2$. Auch die Gleichung $r=2$ ist noch nicht mit einer EABC-Herleitung der komplexen Quantenmechanik gleichzusetzen; sie ist ein Schritt innerhalb von Hardys vollständiger Rekonstruktion \cite{Hardy2001}.
\end{remark}

\section{Diskrete EABC-Signaturen}

Für $n\in\N_{>0}$ sei
\begin{equation}
\Sigma_{EABC}(n)
:=
\bigl(
\Supp_{EABC}(n),\varepsilon(n),\sigma_{EABC}(n)
\bigr)
\in\SEABC,
\end{equation}
wobei
\begin{equation}
\SEABC
:=
\mathcal P(\{E,A,B,C\})\times\FF_2^3\times V_4
\end{equation}
mit der diskreten Topologie versehen wird. Die drei Komponenten bezeichnen den EABC-Familienträger, den ABC-Paritätsvektor und die Restklassensignatur in
\begin{equation}
V_4\cong (\ZZ/12\ZZ)^\times.
\end{equation}
Die Signatur ist eine diskrete arithmetische Invariante; sie ist nicht bereits ein operationaler Quantenzustand.

\section{Das topologische No-Go-Resultat}

\begin{lemma}[Stetige Abbildungen in diskrete Räume \StatusA]
\label{lem:discrete-map}
Es sei $X$ ein wegzusammenhängender topologischer Raum und $D$ ein diskreter Raum. Dann ist jede stetige Abbildung
\[
\Sigma:X\longrightarrow D
\]
konstant.
\end{lemma}

\begin{proof}
Für $x,y\in X$ existiert ein stetiger Weg $\gamma:[0,1]\to X$ mit $\gamma(0)=x$ und $\gamma(1)=y$. Die Komposition $\Sigma\circ\gamma:[0,1]\to D$ ist stetig. Da $[0,1]$ zusammenhängend und $D$ diskret ist, besteht ihr Bild aus genau einem Punkt. Somit gilt $\Sigma(x)=\Sigma(y)$.
\end{proof}

\begin{theorem}[No-Go für eine reine EABC-Signaturdynamik \StatusA]
\label{thm:no-go}
Sei $D$ ein diskreter Raum mit mindestens zwei Punkten. Identifiziert man die reinen Zustände eines Systems unmittelbar mit den Elementen von $D$, so kann zwischen zwei verschiedenen Zuständen kein stetiger Pfad existieren. Hardys Axiom 5 kann daher auf einer nichttrivialen diskreten EABC-Signaturebene nicht erfüllt werden.
\end{theorem}

\begin{proof}
Eine stetige Abbildung $[0,1]\to D$ ist nach Lemma~\ref{lem:discrete-map} konstant. Sie kann deshalb keine zwei verschiedenen Punkte von $D$ verbinden.
\end{proof}

\begin{corollary}[Konstanz auf Hardy-Bahnen \StatusA]
\label{cor:hardy-orbits}
Sei $\Omega$ ein topologischer Zustandsraum, und die reinen Zustände $\Pur(\Omega)$ seien durch kontinuierliche reversible Transformationspfade paarweise verbunden. Dann ist jede stetige diskrete Signatur
\[
\Sigma:\Pur(\Omega)\longrightarrow D
\]
auf $\Pur(\Omega)$ konstant.
\end{corollary}

\begin{remark}[Sektorbegriff]
Aus Korollar~\ref{cor:hardy-orbits} folgt zunächst nur, dass verschiedene EABC-Signaturen nicht verschiedene reine Zustände derselben wegzusammenhängenden Hardy-Bahn unterscheiden können. Neutral gesprochen können sie als \emph{diskrete Sektorlabels}, Faserindizes oder Komponentenlabels dienen. Der stärkere physikalische Begriff eines Superselektionssektors würde zusätzlich Aussagen über zulässige Transformationen, Effekte und kohärente Superpositionen erfordern und wird hier nicht beansprucht.
\end{remark}

\section{Kontinuierliche Zustandsfasern \StatusC}

\begin{hypothesis}[Gefaserter Zustandsraum]
\label{hyp:fibers}
Es existiere ein topologischer Zustandsraum
\begin{equation}
\mathfrak S
=
\bigsqcup_{n\in\N_{>0}}\Omega_n
\end{equation}
mit einer Projektion $\pi:\mathfrak S\to\N_{>0}$ und Fasern $\pi^{-1}(n)=\Omega_n$. Die EABC-Signatur werde nur über den diskreten Basiswert $n$ ausgewertet.
\end{hypothesis}

Dann ergibt sich das kommutative Schema
\begin{equation}
\begin{tikzcd}[row sep=large,column sep=large]
\Pur(\Omega_n)
  \arrow[r,"g_t"]
  \arrow[d,"\pi"']
&
\Pur(\Omega_n)
  \arrow[d,"\pi"]
\\
n
  \arrow[d,"\Sigma_{EABC}"']
&
n
  \arrow[d,"\Sigma_{EABC}"]
\\
\Sigma_{EABC}(n)
  \arrow[r,equal]
&
\Sigma_{EABC}(n).
\end{tikzcd}
\end{equation}
Die kontinuierliche Dynamik verändert den Zustand innerhalb der Faser, nicht aber deren diskretes arithmetisches Label.

\begin{hypothesis}[Transitive Faserwirkung]
\label{hyp:transitive}
Für jedes $n$ wirke die Identitätskomponente $G_n^0$ einer topologischen Gruppe $G_n$ stetig und transitiv auf $\Pur(\Omega_n)$.
\end{hypothesis}

\begin{proposition}[Hinreichende Realisierung der Hardy-Pfade \StatusC]
Unter Hypothese~\ref{hyp:transitive} und der zusätzlichen Voraussetzung, dass jedes Element von $G_n^0$ durch einen Weg mit der Identität verbunden ist, können je zwei reine Zustände derselben Faser durch einen stetigen Pfad reversibler Transformationen verbunden werden.
\end{proposition}

\begin{proof}
Seien $\psi,\phi\in\Pur(\Omega_n)$. Nach Transitivität existiert $g\in G_n^0$ mit $g\psi=\phi$. Nach Voraussetzung gibt es einen Weg $g_t$ in $G_n^0$ von $\id$ nach $g$. Wegen der Stetigkeit der Gruppenwirkung ist $t\mapsto g_t\psi$ ein stetiger Weg von $\psi$ nach $\phi$.
\end{proof}

\begin{remark}
Die Annahme einer Lie-Gruppe ist stärker als Hardys Axiom 5. Ist $G_n^0$ eine zusammenhängende Lie-Gruppe, so ist sie wegzusammenhängend und erfüllt die in der Proposition verwendete Zusatzvoraussetzung. Die Lie-Struktur ist somit eine mögliche hinreichende Realisierung, aber kein Bestandteil des No-Go-Theorems.
\end{remark}

\section{Operatorlift und Komponentengruppe}

\begin{hypothesis}[EABC-Operatorlift \StatusC]
\label{hyp:operator-lift}
Es existiere eine topologische Gruppe $G_{\mathrm{op}}$ und eine kurze exakte Sequenz
\begin{equation}
\begin{tikzcd}
1 \arrow[r]
& G_{\mathrm{cont}} \arrow[r]
& G_{\mathrm{op}} \arrow[r,"q"]
& V_4 \arrow[r]
& 1,
\end{tikzcd}
\end{equation}
wobei
\begin{equation}
G_{\mathrm{cont}}=G_{\mathrm{op}}^0=\ker(q).
\end{equation}
Dann gilt
\begin{equation}
\pi_0(G_{\mathrm{op}})
\cong
G_{\mathrm{op}}/G_{\mathrm{op}}^0
\cong V_4.
\end{equation}
Die vier Zusammenhangskomponenten werden durch $E,A,B,C$ bezeichnet.
\end{hypothesis}

\begin{proposition}[Explizite gespaltene Operatorerweiterung \StatusA]
\label{prop:semidirect}
Sei $G$ eine wegzusammenhängende topologische Gruppe und
\begin{equation}
\alpha:V_4\longrightarrow\Aut_{\mathrm{top}}(G)
\end{equation}
ein Gruppenhomomorphismus in die Gruppe der stetigen topologischen Automorphismen von $G$. Dann ist
\begin{equation}
G_{\mathrm{op}}:=G\rtimes_\alpha V_4
\end{equation}
mit der Produkttopologie eine topologische Gruppe. Es existiert eine gespaltene kurze exakte Sequenz
\begin{equation}
1\longrightarrow G
\longrightarrow G\rtimes_\alpha V_4
\xrightarrow{\ q\ }V_4
\longrightarrow1,
\end{equation}
und
\begin{equation}
(G\rtimes_\alpha V_4)^0=G\times\{E\},
\qquad
\pi_0(G\rtimes_\alpha V_4)\cong V_4.
\end{equation}
\end{proposition}

\begin{proof}
Die Multiplikation ist
\begin{equation}
(g,X)(h,Y)
=
\bigl(g\,\alpha_X(h),X\star Y\bigr).
\end{equation}
Die Projektion $q(g,X)=X$ ist ein surjektiver Gruppenhomomorphismus mit Kern $G\times\{E\}$. Die Abbildung $X\mapsto(e_G,X)$ spaltet $q$. Da $V_4$ diskret ist, sind die Mengen $G\times\{X\}$ offen und abgeschlossen. Weil $G$ wegzusammenhängend ist, sind sie genau die Zusammenhangskomponenten; die Identitätskomponente ist $G\times\{E\}$.
\end{proof}

\begin{remark}[Direktes und semidirektes Produkt]
Für die triviale Wirkung $\alpha_X=\id_G$ erhält man das direkte Produkt $G\times V_4$. Eine nichttriviale Wirkung koppelt die diskreten EABC-Komponenten an die kontinuierliche Operatorgruppe. Proposition~\ref{prop:semidirect} beweist die mathematische Existenz solcher Erweiterungen, nicht ihre physikalische Relevanz.
\end{remark}

\section{Die bedingte \texorpdfstring{$V_4$}{V4}-Fourier-Schnittstelle}

Die Gruppenerweiterung allein erzeugt noch keine spektrale Zerlegung eines Operatorraums. Dafür muss zusätzlich eine lineare Wirkung von $V_4$ auf diesem Raum gegeben sein.

\begin{hypothesis}[Lineare Sektorwirkung \StatusC]
\label{hyp:v4-action}
Es sei $\Oh$ ein komplexer Vektorraum und
\begin{equation}
\rho:V_4\longrightarrow\operatorname{GL}(\Oh)
\end{equation}
eine lineare Darstellung. In einer gespaltenen Operatorerweiterung kann eine solche Wirkung beispielsweise durch Einschränkung einer $G_{\mathrm{op}}$-Darstellung auf die ausgezeichnete Untergruppe $\{e_G\}\times V_4$ entstehen.
\end{hypothesis}

Wir identifizieren
\begin{equation}
V_4\cong\FF_2^2,
\qquad
E=(0,0),\quad A=(1,0),\quad B=(0,1),\quad C=(1,1),
\end{equation}
und wählen die Charaktere
\begin{equation}
\chi_X(Y):=(-1)^{\langle X,Y\rangle},
\qquad X,Y\in\FF_2^2.
\end{equation}
Damit ist eine feste Identifikation $V_4\cong\widehat{V_4}$ gewählt. Die Charaktertafel lautet
\begin{equation}
\begin{array}{c|rrrr}
 & E & A & B & C\\
\hline
\chi_E & 1& 1& 1& 1\\
\chi_A & 1&-1& 1&-1\\
\chi_B & 1& 1&-1&-1\\
\chi_C & 1&-1&-1& 1
\end{array}.
\end{equation}

\begin{theorem}[$V_4$-Fourier-Projektoren \StatusA]
\label{thm:fourier}
Unter Hypothese~\ref{hyp:v4-action} definiert
\begin{equation}
P_X
:=
\frac{1}{4}\sum_{Y\in V_4}\chi_X(Y)\,\rho(Y),
\qquad X\in V_4,
\end{equation}
vier paarweise orthogonale algebraische Projektoren. Es gelten
\begin{equation}
P_XP_Z=\delta_{XZ}P_X,
\qquad
\sum_{X\in V_4}P_X=\id_{\Oh},
\end{equation}
und die direkte Zerlegung
\begin{equation}
\Oh
=
\Oh_E\oplus\Oh_A\oplus\Oh_B\oplus\Oh_C,
\qquad
\Oh_X:=\im(P_X).
\end{equation}
Für $u\in\Oh_X$ gilt äquivalent
\begin{equation}
\rho(Y)u=\chi_X(Y)u
\qquad\text{für alle }Y\in V_4.
\end{equation}
\end{theorem}

\begin{proof}
Die Charakterorthogonalität der endlichen abelschen Gruppe $V_4$ ergibt
\begin{equation}
\frac{1}{4}\sum_{Y\in V_4}\chi_X(Y)\chi_Z(Y)
=\delta_{XZ}.
\end{equation}
Zusammen mit $\rho(Y)\rho(W)=\rho(Y\star W)$ folgen daraus die Projektoridentitäten und die Auflösung der Identität. Die Bildräume sind genau die simultanen Eigenräume der Darstellung.
\end{proof}

\begin{corollary}[Sektorielle Kreuzmultiplikation \StatusA]
\label{cor:sector-product}
Sei $\Oh$ zusätzlich eine komplexe Algebra und wirke $V_4$ durch Algebraautomorphismen. Dann gilt
\begin{equation}
\boxed{
\Oh_X\,\Oh_Z\subseteq\Oh_{X\star Z}
\qquad (X,Z\in V_4).
}
\end{equation}
Die Fourier-Zerlegung macht $\Oh$ damit zu einer $V_4$-graduierten Algebra.
\end{corollary}

\begin{proof}
Für $u\in\Oh_X$, $v\in\Oh_Z$ und $Y\in V_4$ gilt
\begin{align}
\rho(Y)(uv)
&=\rho(Y)(u)\rho(Y)(v)\\
&=\chi_X(Y)\chi_Z(Y)uv\\
&=\chi_{X\star Z}(Y)uv.
\end{align}
Somit liegt $uv$ im Charakterraum $\Oh_{X\star Z}$.
\end{proof}

\begin{remark}[Epistemische Grenze der Fourier-Schnittstelle]
Theorem~\ref{thm:fourier} ist ein formaler Satz über eine \emph{vorgegebene} Darstellung $\rho$. Die Existenz einer physikalisch relevanten Wirkung von $V_4$ auf einer Operatoralgebra folgt weder aus der EABC-Arithmetik noch aus der kurzen exakten Sequenz allein. Diese Existenz und ihre operationale Interpretation gehören zum bedingten Modellrahmen beziehungsweise zu den offenen Forschungsfragen.
\end{remark}

\section{Revidierte Fünf-Ebenen-Architektur}

Die Rollen der verschiedenen Strukturen lassen sich wie folgt ordnen:
\begin{description}[style=nextline,leftmargin=3.2em,labelindent=0em]
  \item[Ebene 0: Skalararithmetik]
  $(\N_{>0},\cdot)$.

  \item[Ebene 1: Diskrete EABC-Signatur]
  $(\SEABC,\odot)$ mit Familienträger, Paritätsvektor und Restklassensignatur.

  \item[Ebene 2: Zustandsfasern]
  $\Omega_n$ und $\Pur(\Omega_n)$.

  \item[Ebene 3: Reversible Faserdynamik]
  $G_n^0\curvearrowright\Pur(\Omega_n)$ als hinreichender Rahmen für kontinuierliche Hardy-Pfade.

  \item[Ebene 4: System- und Operatorkomposition]
  $\Omega_m\boxtimes\Omega_n\to\Omega_{mn}$ sowie
  $\Oh_m\otimes\Oh_n\to\Oh_{mn}$.
\end{description}

Die $V_4$-Fourier-Zerlegung liegt quer zu den Ebenen 1, 3 und 4: Sie übersetzt ein gegebenes diskretes $V_4$-Wirken in vier Operatorsektoren und liefert unter der Automorphismenannahme die Kreuzmultiplikation aus Korollar~\ref{cor:sector-product}.

\begin{remark}[Die Rolle der Komposition \StatusORQ]
Eine tensorielle Operatorenebene kann die Multiplikativität
\begin{equation}
K_{AB}=K_AK_B
\end{equation}
strukturell modellieren. Sie impliziert jedoch allein weder $K=N^r$ noch $K=N^2$ und auch keine komplexe Hilbertraumstruktur. Dafür müssten die vollständigen operationalen Voraussetzungen, insbesondere eine geeignete Zustands-, Effekt-, Subraum- und Kompositionsstruktur, nachgewiesen werden.
\end{remark}

\section{Epistemischer Status}

\subsection*{Quellengestützte Hardy-Aussagen \StatusH}
\begin{itemize}[leftmargin=2.2em]
  \item Innerhalb von Hardys Axiomensystem gilt $K=N^r$ mit $r\in\N_{>0}$.
  \item Axiom 5 schließt $r=1$ aus; das Einfachheitsaxiom wählt danach $r=2$ \cite{Hardy2001}.
\end{itemize}

\subsection*{Formal gesichert \StatusA}
\begin{itemize}[leftmargin=2.2em]
  \item Nichttriviale diskrete EABC-Signaturräume erfüllen die Pfadforderung von Axiom 5 nicht.
  \item Stetige diskrete Signaturen sind auf wegzusammenhängenden Transformationsbahnen konstant.
  \item Direkte und semidirekte Produkte $G\rtimes_\alpha V_4$ liefern gespaltene Gruppenerweiterungen mit Komponentengruppe $V_4$.
  \item Jede komplexe Darstellung von $V_4$ besitzt die Fourier-Projektoren aus Theorem~\ref{thm:fourier}.
  \item Wirkt $V_4$ durch Algebraautomorphismen, so entsteht die $V_4$-graduierte Kreuzmultiplikation $\Oh_X\Oh_Z\subseteq\Oh_{X\star Z}$.
\end{itemize}

\subsection*{Bedingter Modellrahmen \StatusC}
\begin{itemize}[leftmargin=2.2em]
  \item Die Existenz operationaler Zustandsfasern $\Omega_n$ über natürlichen Zahlen.
  \item Transitive kontinuierliche Gruppenwirkungen $G_n^0\curvearrowright\Pur(\Omega_n)$.
  \item Die physikalische Identifikation von $V_4$ mit der Komponentengruppe einer Operatorgruppe.
  \item Eine konkrete lineare oder algebraische $V_4$-Wirkung auf einem physikalisch relevanten Operatorraum.
\end{itemize}

\subsection*{Offene Forschungsfragen \StatusORQ}
\begin{itemize}[leftmargin=2.2em]
  \item Konstruktion einer expliziten Mess- und Effektstruktur;
  \item Nachweis lokaler Tomographie und einer geeigneten Systemkomposition;
  \item Herleitung von $K=N^2$ innerhalb eines EABC-basierten operationalen Modells;
  \item Konstruktion und Klassifikation nichttrivialer Wirkungen
  $\alpha:V_4\to\Aut_{\mathrm{top}}(G_{\mathrm{cont}})$;
  \item Herleitung einer komplexen Hilbertraumdarstellung statt ihrer Postulierung;
  \item operationale Bedeutung der vier Fourier-Sektoren.
\end{itemize}

\begin{theorem}[Präzisiertes Gesamtresultat]
\label{thm:conclusion}
Es gelten folgende Aussagen.
\begin{enumerate}[label=\textup{(\arabic*)},leftmargin=2.5em]
  \item \StatusA{} Die diskrete EABC-Signatur realisiert Hardys fünftes Axiom nicht selbst. Auf jeder wegzusammenhängenden Bahn reiner Zustände ist eine stetige EABC-Signatur konstant.
  \item \StatusA{} Für jede komplexe Darstellung von $V_4$ existiert eine kanonische Charakterzerlegung. Unter einer Wirkung durch Algebraautomorphismen respektiert die Multiplikation die EABC-Regel $X\star Z$.
  \item \StatusC{} Eine Hardy-kompatible Erweiterung kann bedingt durch kontinuierliche Zustandsfasern und transitive reversible Gruppenwirkungen modelliert werden. Ein Operatorlift mit Quotient $V_4$ ist algebraisch konsistent und wird durch Proposition~\ref{prop:semidirect} explizit realisiert.
  \item \StatusORQ{} Aus diesen Strukturen folgen weder $K=N^2$ noch lokale Tomographie, eine komplexe Hilbertraumstruktur oder eine physikalische Identifikation der EABC-Sektoren.
\end{enumerate}
\end{theorem}

\section{Schluss}

Die EABC-Arithmetik liefert ein diskretes Signaturgerüst. Hardys Kontinuitätsaxiom gehört dagegen auf die Ebene kontinuierlicher Zustandsbahnen. Die mathematisch tragfähige Brücke besteht deshalb nicht in einer Identifikation beider Ebenen, sondern in einer Faser- und Erweiterungsarchitektur:
\begin{equation}
\N_{>0}
\xrightarrow{\ \Sigma_{EABC}\ }
\SEABC,
\qquad
1\longrightarrow G_{\mathrm{cont}}
\longrightarrow G_{\mathrm{op}}
\longrightarrow V_4
\longrightarrow1.
\end{equation}
Eine vorgegebene $V_4$-Wirkung auf einer Operatoralgebra erzeugt zusätzlich eine Fourier-Zerlegung und eine $V_4$-graduierte Kreuzmultiplikation. Das No-Go-Theorem und die Fourier-Sätze sind formal; die operationale und physikalische Interpretation des Operatorlifts bleibt ein Forschungsprogramm.

\begin{thebibliography}{99}
\bibitem{Hardy2001}
L. Hardy,
\emph{Quantum Theory From Five Reasonable Axioms},
arXiv:quant-ph/0101012v4, 2001.

\bibitem{Serre1977}
J.-P. Serre,
\emph{Linear Representations of Finite Groups},
Springer, 1977.

\bibitem{Hoffbauer2026}
T. Hoffbauer,
\emph{Skalare Projektionen, Faserquotienten und EABC-Paritätsebenen in Modellräumen},
Preprint, August 2026.
\end{thebibliography}

\end{document}
